\documentclass[11pt,a4paper]{article}
\usepackage[margin=0.95in]{geometry}
\usepackage[T1]{fontenc}
\usepackage{lmodern}
\usepackage{amsmath,amssymb}
\usepackage{xcolor}
\usepackage{fancyhdr}
\usepackage{hyperref}
\setlength{\parskip}{0.5em}
\setlength{\parindent}{0pt}
\setlength{\headheight}{15pt}
\linespread{1.08}
\pagestyle{fancy}
\fancyhf{}
\lhead{CST444 Soft Computing - Module 1}
\rhead{PDF-Source Locked Guide}
\cfoot{\thepage}
\hypersetup{colorlinks=true,linkcolor=blue,urlcolor=blue}

\newcommand{\examtip}[1]{%
\vspace{0.2em}
\noindent\fcolorbox{black}{gray!12}{\parbox{0.97\linewidth}{\textbf{Exam Tip:} #1}}%
\vspace{0.2em}
}

\begin{document}

\section*{Source Lock Declaration}
This Module 1 guide is rewritten using only \textbf{SC - M1 -Ktunotes.in.pdf}. Question coverage follows your OCR question set, but answer points are sourced from that PDF content and its worked-problem sections.

\section*{Module 1 Topic Map from Questions}
\textbf{Frequently repeated concepts}
\begin{itemize}
\item Soft computing vs hard computing.
\item Biological neuron and artificial neuron comparison.
\item Net input, bias, and activation function role.
\item McCulloch-Pitts neuron and logic implementation.
\item Hebb rule and Hebb logical designs.
\item Linear separability and XOR non-separability.
\end{itemize}

\textbf{High-mark concepts}
\begin{itemize}
\item Linear separability with XOR justification.
\item AND/ANDNOT using MP neuron with threshold logic.
\item Hebb training algorithm with bipolar coding.
\item Activation-function based network outputs.
\end{itemize}

\textbf{Must-study-first order}
\begin{enumerate}
\item Soft computing idea and characteristics.
\item Biological neuron parts and behavior.
\item Artificial neuron model and bias effect.
\item Activation functions and output ranges.
\item MP neuron logic realization.
\item Hebb learning and XOR limitation.
\end{enumerate}

\section*{Core Theory (From Source PDF)}

\subsection*{1) Soft Computing and Hard Computing}
\textbf{From source notes}
\begin{itemize}
\item Soft computing was proposed by Lotfi A. Zadeh.
\item It exploits tolerance for imprecision and uncertainty.
\item It aims at tractability, robustness, and low solution cost.
\item The role model for soft computing is the human mind.
\item Main techniques listed are ANN, fuzzy logic, and GA.
\item Hard computing emphasizes precise result and formal model.
\end{itemize}

\subsection*{2) Biological Neuron and ANN Mapping}
\begin{itemize}
\item Biological neuron parts listed: soma, dendrites, axon, synapse.
\item Dendrites receive signals; axon carries impulses.
\item Firing behavior depends on crossing threshold.
\item ANN is described as inspired by biological nervous system.
\item ANN uses interconnected neurons with weighted links.
\item Weights contain input signal information.
\end{itemize}

\subsection*{3) Net Input and Bias}
\begin{itemize}
\item Source notation uses \textbf{Yin} for net input and \textbf{Y} for output.
\item Bias is included by adding an input term ($x_0=1$ style statement in notes).
\item Positive bias increases net input.
\item Negative bias decreases net input.
\item Net-input questions are repeatedly asked in module problems.
\item Output is obtained after applying activation over net input.
\end{itemize}

\subsection*{4) Activation Functions}
\begin{itemize}
\item Activation function is applied over net input to obtain output.
\item Linear activation has limited complexity handling.
\item Nonlinear activation is needed for complex mappings.
\item Sigmoid is non-linear and bounded.
\item Binary sigmoid output is between 0 and 1.
\item Bipolar sigmoid output is between -1 and 1.
\end{itemize}

\subsection*{5) McCulloch-Pitts Neuron}
\begin{itemize}
\item MP neuron is an early binary computational neuron model.
\item Inputs are binary and output is threshold-based.
\item Links can be excitatory or inhibitory.
\item Firing is based on threshold condition.
\item Source notes state simple MP neurons can realize logical operations.
\item AND and ANDNOT are included in source problem sets.
\end{itemize}

\subsection*{6) Hebb Learning and Linear Separability}
\begin{itemize}
\item Hebb learning idea: synchronous activation strengthens connection.
\item Source states Hebb rule is more suited to bipolar data.
\item Hebb sections include AND, OR, XOR, and pattern-classification examples.
\item Linear separability is explained with one separating line in 2D.
\item XOR is explicitly given as linearly inseparable.
\item Notes state XOR cannot be separated with a single line.
\end{itemize}

\examtip{Use source phrase in exams: ``XOR cannot be separated with a single line.'' This is a direct scoring statement from your notes.}

\section*{Numericals and Derivations (Source-Only)}

\subsection*{N1) Activation numerical from source worked example}
\textbf{Given in source problem}
\begin{itemize}
\item Net input to neuron $Y$ is shown as 0.93.
\item Binary sigmoidal activation is asked.
\end{itemize}

\textbf{Result reported in source}
\begin{itemize}
\item Output of neuron $Y$ using binary sigmoidal activation is 0.72.
\end{itemize}

\textbf{Exam format writeup}
\begin{itemize}
\item Given: $Y_{in}=0.93$.
\item Formula used: binary sigmoid function (as source activation section).
\item Final answer (source result): $Y=0.72$.
\end{itemize}

\subsection*{N2) Net-input diagram problems (source method)}
\textbf{Given:} network diagram with inputs, weights, and bias.

\textbf{Find:} net input (Yin), and output if activation is specified.

\textbf{Source-based method}
\begin{itemize}
\item Read all weighted links from diagram.
\item Include bias influence (increase/decrease Yin based on sign).
\item Compute aggregated net input Yin.
\item Apply requested activation function.
\item Report output in required range.
\end{itemize}

\subsection*{N3) Hebb pattern classification (I vs O) from source section}
\textbf{Given in notes}
\begin{itemize}
\item 3x3 pattern matrix.
\item ``+'' indicates 1 and blank indicates -1.
\item I has target 1 and O has target -1.
\end{itemize}

\textbf{Source-based solution flow}
\begin{itemize}
\item Initialize weights and bias to zero.
\item Present each bipolar pattern with target.
\item Update weights using Hebb learning statement.
\item Test final weights on I and O.
\item Verify sign of output against target class.
\end{itemize}

\section*{Solved Questions for Module 1}

\subsection*{Part A / 3-Mark Questions (Each Answer Has 6 Scoring Points)}

\textbf{Q1. Differentiate soft computing and hard computing.}
\begin{enumerate}
\item Soft computing handles imprecision and uncertainty.
\item Hard computing requires precise formulation.
\item Soft computing targets robustness and low solution cost.
\item Hard computing gives precise guaranteed output.
\item Soft computing is adaptive and human-mind inspired.
\item ANN, fuzzy logic, GA are listed soft-computing techniques.
\end{enumerate}

\textbf{Q2. Role of activation function; explain any two.}
\begin{enumerate}
\item Activation is applied over net input to produce output.
\item It decides whether neuron response should be active.
\item Linear activation has limited expressive ability.
\item Nonlinear functions support complex mapping.
\item Binary sigmoid outputs lie between 0 and 1.
\item Bipolar sigmoid outputs lie between -1 and 1.
\end{enumerate}

\textbf{Q3. Draw simple artificial neuron and discuss net input.}
\begin{enumerate}
\item Inputs are connected through weighted links.
\item Yin is treated as net input in source notation.
\item Output Y is obtained after activation operation.
\item Weights represent signal strength.
\item Bias affects Yin magnitude and direction.
\item Net-input calculation is central to ANN response.
\end{enumerate}

\textbf{Q4. Compare biological and artificial neuron.}
\begin{enumerate}
\item Biological neuron has soma, dendrites, axon, synapse.
\item ANN neuron is a simplified processing unit.
\item Dendritic reception maps to ANN input stage.
\item Synaptic behavior maps to ANN weighted connections.
\item Biological firing is threshold dependent.
\item ANN output is activation-function dependent.
\end{enumerate}

\textbf{Q5. Explain ANN architectures based on connection.}
\begin{enumerate}
\item Single layer feedforward network.
\item Multilayer feedforward network.
\item Single node with its own feedback.
\item Single layer recurrent network.
\item Multilayer recurrent network.
\item Source also groups as feedforward/recurrent categories.
\end{enumerate}

\textbf{Q6. ``Hebb rule is more suited for bipolar data.'' Justify.}
\begin{enumerate}
\item This point is explicitly stated in source notes.
\item Hebb rule ties weight change to input-learning-signal relation.
\item Bipolar coding supports positive/negative correlation learning.
\item Source examples of Hebb logic are in bipolar form.
\item XOR/AND/OR Hebb examples all follow bipolar style.
\item Therefore bipolar representation is preferred in source treatment.
\end{enumerate}

\textbf{Q7. Explain Hebb training algorithm.}
\begin{enumerate}
\item Start with initial weights and bias (source problems use zero start).
\item Present one input-target pattern.
\item Apply Hebb weight-change relation.
\item Carry updated weights to next pattern.
\item Continue until all training patterns are presented.
\item Use final weights for testing/classification.
\end{enumerate}

\textbf{Q8. Types of learning methods in neural networks.}
\begin{enumerate}
\item Supervised learning is target-guided.
\item Unsupervised learning works without teacher targets.
\item Reinforcement learning uses reward feedback.
\item Source discusses adaptive learning as ANN property.
\item Learning modifies synaptic weights over time.
\item Updated weights determine future classification quality.
\end{enumerate}

\textbf{Q9. Implement AND using MP neuron (binary).}
\begin{enumerate}
\item Source includes dedicated MP-AND problem sequence.
\item Use binary inputs as MP model requires.
\item Compute summed input for each truth-table row.
\item Apply threshold decision in each case.
\item Output is high only for AND-true case.
\item Thus MP neuron realizes AND logic.
\end{enumerate}

\textbf{Q10. Implement ANDNOT using MP neuron.}
\begin{enumerate}
\item Source includes explicit MP-ANDNOT problem sequence.
\item Use binary inputs and threshold activation.
\item Include inhibitory interpretation for NOT part.
\item Evaluate truth table through threshold condition.
\item Only ANDNOT-true case produces high output.
\item Therefore MP neuron models ANDNOT function.
\end{enumerate}

\textbf{Q11. Net input calculation question.}
\begin{enumerate}
\item Read all inputs and weight values from diagram.
\item Include bias term and its sign effect.
\item Aggregate to obtain Yin.
\item Check if activation function is additionally requested.
\item Apply requested activation to Yin.
\item Report final output with proper range.
\end{enumerate}

\textbf{Q12. Sigmoid-output question.}
\begin{enumerate}
\item Source includes binary and bipolar sigmoid examples.
\item First compute Yin from the given network.
\item For binary sigmoid, output range is [0,1].
\item For bipolar sigmoid, output range is [-1,1].
\item Source numeric example: Yin 0.93 gives binary output 0.72.
\item Write final answer with requested activation type.
\end{enumerate}

\subsection*{Part B / 14-Mark Questions (Each Answer Has 14 Scoring Points)}

\textbf{Q1. Explain linear separability and justify XOR non-linearity; include AND/Hebb context.}
\begin{enumerate}
\item Linear separability means one line separates two classes in 2D.
\item Notes describe class points occupying opposite half-spaces.
\item A single-layer classifier depends on such separability.
\item AND-like class setting is separable in standard 2D view.
\item Source includes Hebb logic examples for separable cases.
\item XOR is listed as classic non-separable pattern.
\item XOR class points are diagonally distributed.
\item No single line can split XOR classes correctly.
\item Therefore single linear boundary is insufficient for XOR.
\item Source explicitly says XOR cannot be separated by one line.
\item This explains failure of simple linear classification.
\item More complex partitioning is required for XOR.
\item In exam, add a 2D statement about opposite diagonals.
\item Conclusion: separability criterion explains both AND success and XOR failure.
\end{enumerate}

\textbf{Q2. Implement ANDNOT using MP neuron with architecture and threshold condition.}
\begin{enumerate}
\item MP neuron uses binary input and threshold output.
\item Source states links may be excitatory/inhibitory.
\item Build two-input binary ANDNOT truth table.
\item Compute summed input per row.
\item Compare each sum with fixed threshold.
\item If threshold is satisfied, neuron fires.
\item Otherwise output remains zero.
\item Input-1 supports the AND term.
\item Input-2 supports NOT behavior via inhibition.
\item Verify four input cases one by one.
\item Exactly one case should produce output 1.
\item Remaining three cases produce output 0.
\item This matches ANDNOT functional definition.
\item Conclusion: MP threshold architecture realizes ANDNOT directly.
\end{enumerate}

\textbf{Q3. Using Hebb rule, derive classifying weights for I and O pattern and explain testing.}
\begin{enumerate}
\item Source problem defines 3x3 pattern classification setup.
\item Convert symbols to bipolar values (+ as 1, blank as -1).
\item Assign targets: I as +1 and O as -1.
\item Initialize all weights and bias to zero.
\item Present first pattern and apply Hebb update relation.
\item Carry updated weights into next pattern.
\item Continue until all patterns are presented.
\item Maintain bipolar consistency through entire process.
\item Use final learned weights for testing stage.
\item Test pattern I and check output sign.
\item Test pattern O and check output sign.
\item Compare outputs with assigned targets.
\item Report classification correctness.
\item Conclusion: Hebb network performs source-defined class separation for the I/O patterns.
\end{enumerate}

\textbf{Q4. Solve activation-function output question (binary and bipolar sigmoid) in source style.}
\begin{enumerate}
\item Identify given network and calculate net input Yin.
\item Include bias contribution exactly as shown.
\item If binary sigmoid is asked, use source binary range [0,1].
\item If bipolar sigmoid is asked, use source bipolar range [-1,1].
\item Apply activation over computed Yin.
\item Present binary output.
\item Present bipolar output.
\item Compare ranges and behavior.
\item State why different activations give different values.
\item Mention non-linearity role from source activation section.
\item Keep answer in Given-Find-Method-Final structure.
\item For source example: Yin 0.93 gives binary output 0.72.
\item If diagram data is figure-only, retain method and source reporting style.
\item Conclusion: activation choice controls output scale and response profile.
\end{enumerate}

\section*{Formula/Method Sheet (Source-Safe)}
\begin{itemize}
\item Net input notation in notes: Yin.
\item Output notation in notes: Y.
\item MP rule: output 1 if threshold is satisfied, else 0.
\item Binary sigmoid output range: 0 to 1.
\item Bipolar sigmoid output range: -1 to 1.
\item Hebb statement: weight increase is proportional to input and learning signal product.
\item Linear separability: one line can separate two classes.
\item XOR: cannot be separated with a single line.
\end{itemize}

\section*{One-Page Quick Revision}
\begin{itemize}
\item Soft computing is uncertainty-tolerant and adaptive.
\item Hard computing is precise and model-centric.
\item ANN mirrors biological-network ideas with weighted links.
\item Bias shifts net input up or down.
\item Activation maps net input to final output.
\item MP neuron is binary threshold logic.
\item Hebb uses local connection strengthening behavior.
\item XOR is the standard non-separable logic case.
\end{itemize}

\examtip{When writing fast in exam, anchor each answer with one direct phrase from your notes and then list scoring points.}

\end{document}
